\chapter{INTRODUCTION}

L’analyse de longs flux vidéo constitue un défi technique significatif en raison du volume important de données à traiter et des coûts élevés en ressources informatiques associés aux approches llms multimodales actuellement disponibles. Ces modèles exigent une puissance de calcul élevée, entraînant des coûts substantiels lorsqu’exécutés sur des serveurs infonuagiques via des API ou une grande puissance de calcul lorsqu’exécutés sur des plateformes embarquées, ce qui limite leur utilisation pratique dans des contextes où les ressources informatiques ou financières sont restreintes. 

Pour réaliser cette tâche avec peu de ressources, l'approche multimodale conventionnelle n’est pas adaptée. L’architecture est plus sophistiquée qu’un LLM textuel, car elle doit utiliser des encordeurs spécialisés aux différents médias et les fusionner en utilisant la fusion profonde via de l’attention croisée et de la fusion précoce \cite{wadekar2024_evolution_of_multimodal_model_architectures}. Cette complexité mène à des entrainements plus coûteux qui requiert significativement plus de données et de puissance de calcul qu’un LLM textuel et à une structure plus volumineuse qui mène à un plus grand nombre de paramètres \cite{liang2025_mixture_of_transformers}. 

Une alternative émergente consiste à exploiter les avancées récentes obtenues avec les grands modèles de langage (LLM), qui atteignent maintenant un niveau remarquable en matière de compréhension et d'analyse sémantique de textes \cite{zheng2023gptfathom}. Cette recherche propose de convertir et d’ordonner l'intégralité des données audiovisuelles issues d'un robot en une représentation purement textuelle, pour que le traitement analytique soit confié exclusivement à un modèle linguistique, réduisant ainsi la dépendance aux méthodes multimodales complexes et coûteuses en ressources \cite{liang2025_mixture_of_transformers}.

Cette solution soulève quand-même plusieurs questions techniques et scientifiques importantes. Il faut vérifier si la représentation textuelle de la vidéo peut être assez fidèle et exhaustive pour permettre une analyse complète de la vidéo par le texte. Une description textuelle complète nécessite non seulement de retranscrire fidèlement les dialogues ou les voix hors champ, mais pourrait également bénéficier de capter efficacement l’essentiel des informations visuelles et non verbales telles que les expressions faciales, le langage corporel, les éléments contextuels et les transitions scénaristiques. Il faut également déterminer si les grands modèles de langage textuel ont une compréhension assez avancée pour identifier correctement les moments clés d'une vidéo avec une précision comparable à celle obtenue par un humain en utilisant la représentation textuelle du vidéo. 

L'objectif scientifique principal est de démontrer la faisabilité technique d'un processus entièrement textuel de détection d'événements, adapté aux contraintes spécifiques de l’analyse de long flux vidéos avec des ressources matérielles ou monétaires limitées. Le succès de cette approche mènerait à l’implémentation d’analyse vidéo automatisée sur de long flux vidéos dans des applications qui ne pouvait jusqu’à présent pas justifier les coûts computationnel et monétaires associés. Comme analyser des flux de caméras de surveillance résidentielle ou automatiser l’assurance qualité d’un centre d’appel en vidéo-conférence en recueillant les moments problématiques.

Sur le web, plusieurs plateformes alimentées par l’IA ont émergé pour simplifier et accélérer le processus de création de clips courts. Ces plateformes offrent aux créateurs de contenu, aux équipes marketing et aux utilisateurs sans compétences la possibilité de générer des clips attractifs à partir des moments saillants issus de vidéos longues. Même si elles ne sont pas adaptées pour la recherche académique, elles comblent un besoin très semblable à ceux des chercheurs qui désirent extraire les moments forts d’un long flux vidéo avec des ressources limitées. Elles doivent générer un très grand nombre de clips pour un très faible coût monétaire dû à la forte compétition entre plateformes. Ces solutions purement logicielles et automatisées gardent leurs coûts bas en utilisant des modèles de langage optimisés et des techniques de traitement vidéo à faible puissance de calcul. Ce choix technologique est lié aux mêmes contraintes rencontrées par les chercheurs qui désirent analyser le flux vidéo issu des capteurs embarqués sur des appareils disposant de ressources computationnelles limitées comme des robots ou des ordinateurs portables. De manière similaire aux plateformes web devant minimiser leurs coûts de calcul pour offrir des solutions abordables aux consommateurs, les chercheurs pourraient bénéficier d'approches légères d’analyse vidéo pour conserver des performances acceptables avec des ressources computationnelles limitées. Ainsi, étudier les techniques employées par ces plateformes web offre des pistes utiles pour développer des méthodes d'analyse vidéo adaptées aux systèmes embarqués avec ressources limitées.

Le chapitre 2 présente l’état de l’art de la détection de moments forts dans de long flux vidéos et analyse les solutions webs existante. 
Le chapitre 3 explique la problématique en détails. 
Le chapitre 4 propose une approche pour régler la problématique et montre une méthodologie et un échéancier possible pour mener à terme la maitrise.





